\section{Methodology}

We wrote a C program, \texttt{batch\_server}, with a known performance
bug. It runs in rounds, about 11 per second: a coordinator thread hands
work to four worker threads, each computing for about 5\,ms, waits until
all of them finish, then idles for 50\,ms. One worker, \texttt{worker-2},
also commits its results to a shared cache guarded by a mutex,
\texttt{cache\_lock}, and a \texttt{flusher} thread writes the cache back
while holding that lock, simulated by a 40\,ms sleep. All six threads
share one CPU. Holding a lock across slow I/O is a common concurrency
mistake: \texttt{worker-2} waits 24.7\,ms per round for the lock, so a
round takes 41.8\,ms instead of the roughly 20\,ms that four 5\,ms
computations on one CPU need. The program logs each round's timings,
which provide the ground truth.

We traced the program system-wide with LTTng kernel tracing (system calls,
scheduling, timers and interrupts) for 2, 10 and 30 minutes, giving traces
of 5.9, 29.3 and 86.9 million events (1{,}305, 6{,}528 and 19{,}580
rounds). Activity from the rest of the machine stays in the traces, as it
would in practice. Three systems, all using GPT-5.6 through the OpenAI
API, answer the same questions about each trace.

\subsection{Baseline 1: LLM Only}

The first baseline tests whether a capable LLM can answer from the trace
itself, with no tools. Each trace is printed as text with
\texttt{babeltrace2}, one line per event. Even the 2-minute trace has
5.9 million events, about 380 million tokens, so the model receives a
\emph{slice} of at most 200{,}000 tokens, cut in one of three ways.
\textbf{(A)~Head:} the first 200k tokens, which is what a user gets by
pasting the trace into a chat; it covers only the first 4.7--5.7\,ms
(about 2{,}800 events) and ends before the program starts.
\textbf{(B)~Middle:} 200k tokens of consecutive events from the trace's
midpoint, about 3{,}200 events spanning 96--125\,ms, roughly one round.
\textbf{(C)~Sampled:} 20 windows of 10k tokens spaced evenly across the
trace, each labelled with its time range, trading depth for coverage.

A header gives the full trace's length and says most events are omitted.
The system prompt asks the model to act as a Linux tracing expert, answer
only from the trace text, name threads, thread IDs, CPUs and durations,
and say when the text is not enough. Each question is a separate API call,
so no answer can draw on another, and questions name threads by name only.
The model never sees the ground-truth files. With 13 questions, 3 slices
and 3 traces, the baseline gives 117 answers, each graded once.

\subsection{Baseline 2: LLM with RAG}

The second baseline adds retrieval to the first and changes nothing else.
Its knowledge base holds 38 documents: 15 on Linux kernel tracing concepts,
such as futex waits, thread states and wake-ups; 12 on Trace Compass
analyses; 10 on Django web applications from an unrelated workload, which
act as distractors; and a description of \texttt{batch\_server}, standing
in for the internal documentation an organization keeps for its own
software. It names the threads and \texttt{cache\_lock} and states the
design, including that the flusher sleeps for 40\,ms while holding the
lock, so it points to the likely root cause, as such documentation would in
practice. It contains no measurements of any run, so it cannot supply the
measured answers, such as how long \texttt{worker-2} waits per round or how
much faster rounds would be after a fix; those must come from the trace.
\sysname uses the same 38 documents, so the documentation is not what
separates the two systems.

Documents are embedded with \texttt{bge-small-en-v1.5} and stored in
ChromaDB. The question is the search query: the six nearest documents are
kept, and the top three are expanded to their full text. This is
\sysname's retrieval layer, except that \sysname's query also draws on its
trace findings. The retrieved documents are fixed per question; the program
description was retrieved for 10 of the 13 questions, ranked first for 6.
The system prompt adds one instruction: take facts about the system from
the trace text, and use the documents only to interpret it. An answer that
is correct only by repeating the program description, with no support in
the trace slice, scores half.

\subsection{\sysname}

\sysname is a Trace Compass plug-in. A deterministic planner maps each
question to a target (a thread, a process or the whole trace) and an
intent, such as explaining where a thread's time went or how a process's
threads hold each other up. Java analyzers then compute findings from the
whole trace; RAG adds documentation as further context, such as what a
futex wait is or what the program was designed to do; and an LLM, the same
GPT-5.6 model as in both baselines, combines the findings and documents
into the answer.

\textbf{Java analyzers.} Three analyzers, built on Trace Compass's kernel
state system and execution graph, cover the three targets:
\begin{itemize}
  \item The \emph{trace explainer} summarizes the trace: which events were
        recorded, which processes ran and used the CPU, and which are the
        workload rather than background activity or the tracer.
  \item The \emph{process explainer} splits each thread's time into
        running, waiting for CPU and blocked, follows the execution graph's
        wake-up edges to find which thread ended each wait, and names the
        bottleneck: the thread the main thread waits on, with its wait
        chain to the root (here main thread $\leftarrow$
        \texttt{worker-2} $\leftarrow$ \texttt{flusher} $\leftarrow$
        timer). Over every round it measures when each thread finishes, how
        long the lock holder holds the lock, inferred from the waits and
        wake-ups around it, and how long rounds would take without the
        bottleneck's wait.
  \item The \emph{thread explainer} splits one thread's time by state and
        explains each large share: for blocked time, the system call and
        the critical path to the thread that ended the wait; for CPU waits,
        which threads held the CPU; for running time, user and system time.
        It also adds the thread's part in its process's recurring wait
        chain, so that a one-off wait, such as the startup pause, is not
        mistaken for the recurring cause.
\end{itemize}
Each result is a \emph{finding}: a JSON object with a one-sentence claim
and its numbers, the time window, the evidence, a confidence level and
caveats. Inferences a kernel trace cannot observe directly, such as lock
ownership, are marked medium confidence. The LLM's rules require every
statement about the trace to come from a finding, with its numbers, and a
missing finding to be stated rather than guessed. An auxiliary statistical
module, TMLL, adds system-call summaries; we leave its evaluation to
future work.

\textbf{Protocol and limitation.} \sysname currently identifies the thread
a question asks about by its ID, as in ``Why is worker-2 (tid 164857)
slow?''. A name resolver handles process names and uniquely named threads,
but threads that share a name, as in a thread pool, still need the ID, so
our thread questions carried it; the baselines' questions did not.
Removing this dependence is future work. For each trace, the questions
were asked in two batches (Q1--Q10 and Q11--Q13), each one conversation in
which later answers can see earlier ones, giving 39 answers.
